\documentclass[letterpaper,10pt]{article}

\usepackage[T1]{fontenc}
\usepackage{charter}
\usepackage[margin=0.75in,top=0.6in,bottom=0.6in]{geometry}
\usepackage{enumitem}
\usepackage{xcolor}
\usepackage[hidelinks]{hyperref}

\pagestyle{empty}
\setlength{\parindent}{0pt}
\raggedright

% Section title over a full-width rule, as in the earlier versions.
\newcommand{\cvsection}[1]{%
  \vspace{9pt}{\large\bfseries #1}\par\vspace{-5pt}%
  \rule{\linewidth}{0.4pt}\par\vspace{1pt}}

% Organisation | role on the left, period flush right.
\newcommand{\cvheading}[3]{\textbf{#1} $|$ #2\hfill #3\par}

% Achievements are indented lines without bullets.
\newenvironment{cvitems}
  {\begin{itemize}[leftmargin=0.16in,label={},itemsep=1.5pt,topsep=2pt,parsep=0pt]}
  {\end{itemize}}

\newcommand{\cvstack}[2]{\vspace{2pt}\textbf{#1:} #2\par}

\begin{document}

\begin{center}
  {\fontsize{25}{29}\selectfont Guilherme Peres}\par\vspace{3pt}
  {\large Desenvolvedor Full Stack}\par\vspace{5pt}
  Santos - São Paulo $|$
  \href{mailto:guilhermeperes.dev@gmail.com}{guilhermeperes.dev@gmail.com} $|$
  (13) 98862-7400\par\vspace{2pt}
  % Full URLs stay in the text because ATS parsers often drop link annotations.
  {\small\color{black!60}%
  \href{https://guilhermeperes.dev/pt}{guilhermeperes.dev} $|$
  \href{https://linkedin.com/in/guilhermeperesdev}{linkedin.com/in/guilhermeperesdev} $|$
  \href{https://github.com/guilhermedkdk}{github.com/guilhermedkdk}}
\end{center}

\cvsection{Perfil}
Desenvolvedor Full Stack com 2 anos de experiência em plataformas financeiras e sistemas ERP.
Trabalho com Java e Node.js no backend, React e Angular no frontend, AWS e IA generativa
aplicada ao negócio.

\cvsection{Experiência Profissional}
\cvheading{Anka Tech}{Desenvolvedor Full Stack}{Set 2025 - Atual}
\textit{Software house especializada em escritórios de assessoria de investimentos.}\par
\begin{cvitems}
  \item Criei recursos de IA generativa com a OpenAI, como uma consulta em linguagem natural à
        base de clientes que transforma perguntas em filtros validados, sem enviar dados
        pessoais ao modelo.
  \item Desenvolvi o motor de comissões mensais com regras configuráveis por assessor, eliminando
        os erros manuais.
  \item Construí pipelines ETL diários com a API do BTG em AWS Lambda, consolidando 3 fontes de dados.
  \item Implementei OAuth2 com SSO e controle de acesso RBAC, isolando as informações por
        assessor e por equipe.
  \item Desenvolvi em Spring Boot e Angular o módulo de carteiras, em que os assessores acompanham
        posições e rentabilidade dos clientes em tempo real, com cotações via WebSocket abaixo de 200 ms.
\end{cvitems}
\cvstack{Stack}{React, Next.js, Node.js, TypeScript, Java, Spring Boot, Angular, PostgreSQL, AWS}

\vspace{7pt}
\cvheading{Audax Electronics}{Desenvolvedor Web Full Stack Jr}{Out 2024 - Set 2025}
\textit{Indústria que projeta e fabrica soluções em LED sob medida para grandes marcas.}\par
\begin{cvitems}
  \item Desenvolvi o ERP interno da empresa, com 10 módulos em microsserviços (produção,
        compras, logística, controladoria, qualidade e CRM).
  \item Criei o pipeline que sincroniza 25 tabelas do TOTVS Protheus com fila e agendamento
        automático, eliminando 4 horas diárias de trabalho manual do time de TI.
  \item Acelerei em mais de 15x as telas mais acessadas do ERP, de $\sim$7 s para menos de 400 ms,
        otimizando consultas e índices no PostgreSQL.
  \item Reduzi a emissão de relatórios operacionais de $\sim$30 min para menos de 5 s com geração
        dinâmica de PDF.
\end{cvitems}
\cvstack{Stack}{Java, Spring Boot, Angular, TypeScript, JPA Hibernate, PostgreSQL}

\vspace{7pt}
\cvheading{SOC - Software de Saúde e Segurança do Trabalho}{Estágio em Suporte}{Out 2023 - Jun 2024}
\begin{cvitems}
  \item Prestei suporte técnico por chat e telefone aos usuários, com registro e escalonamento de chamados.
\end{cvitems}

\cvsection{Projetos}
\textbf{RPGForge AI}\hfill\href{https://rpgforge-ai.vercel.app}{rpgforge-ai.vercel.app}\par
\begin{cvitems}
  \item Plataforma de fichas de RPG com IA, com geração ancorada por RAG em pgvector, saída
        validada por schema Zod, motor de regras próprio, controle de custo por token e 271
        testes automatizados.
\end{cvitems}
\cvstack{Tecnologias}{TypeScript, Next.js, NestJS, Prisma, PostgreSQL, pgvector, OpenAI, Docker}

\cvsection{Habilidades}
\textbf{Linguagens:} TypeScript, JavaScript, Java, Python, SQL\par
\textbf{Backend:} Node.js, NestJS, Spring Boot, FastAPI, REST, WebSocket, Microsserviços\par
\textbf{Frontend:} React, Next.js, Angular, Tailwind CSS\par
\textbf{Dados e Cloud:} PostgreSQL, MongoDB, MySQL, AWS (Lambda, API Gateway, S3, RDS, ECS, CloudWatch), Docker\par
\textbf{IA:} OpenAI API, RAG, pgvector, Embeddings, Structured Outputs\par

\cvsection{Formação e Idiomas}
\textbf{Universidade Católica de Santos}, Ciência da Computação\hfill Jan 2021 - Jun 2025\par
\textbf{Cambridge Assessment English}, Certificado de Inglês CEFR B1\par

\end{document}
